%==============================================================
%  PHẦN ĐẦU: Lời cam đoan, Lời cảm ơn, Danh mục từ viết tắt
%==============================================================

\chapter*{Lời cam đoan}
\addcontentsline{toc}{chapter}{Lời cam đoan}
\markboth{Lời cam đoan}{Lời cam đoan}

Tôi xin cam đoan đây là công trình nghiên cứu của riêng tôi, được thực hiện dưới
sự hướng dẫn của \thesisAdvisor. Các kết quả lý thuyết được trình bày trong luận
văn hoặc do tôi tự chứng minh, hoặc được trích dẫn đầy đủ từ tài liệu gốc. Toàn
bộ số liệu thực nghiệm được tạo ra bằng mã nguồn công bố kèm luận văn
(Phụ lục~\ref{app:code}) và có thể kiểm chứng lại; các giới hạn về tính tái lập
của những số liệu ấy được nêu thẳng ở Mục~\ref{sec:hanche}. Luận văn không sao
chép từ bất kỳ công trình nào khác mà không ghi rõ nguồn.

\vspace{1.5cm}
\begin{flushright}
\begin{minipage}{7cm}
\centering
TP. Hồ Chí Minh, \ngaythangnam{5}{9}{2026}\\[0.3cm]
\textit{Sinh viên thực hiện}\\[2cm]
\thesisAuthor
\end{minipage}
\end{flushright}

%--------------------------------------------------------------
\chapter*{Lời cảm ơn}
\addcontentsline{toc}{chapter}{Lời cảm ơn}
\markboth{Lời cảm ơn}{Lời cảm ơn}

\TODO{Phần này cần tác giả tự viết. Thông thường gồm: lời cảm ơn giảng viên
hướng dẫn, quý thầy cô Khoa Toán--Tin học, gia đình và bạn bè. Xoá lệnh
\texttt{\textbackslash TODO} này sau khi viết xong.}

%--------------------------------------------------------------
\chapter*{Danh mục từ viết tắt}
\addcontentsline{toc}{chapter}{Danh mục từ viết tắt}
\markboth{Danh mục từ viết tắt}{Danh mục từ viết tắt}

\begin{longtable}{@{}l p{5.3cm} p{6.0cm}@{}}
\toprule
\textbf{Viết tắt} & \textbf{Nguyên văn} & \textbf{Nghĩa dùng trong luận văn} \\
\midrule
\endfirsthead
\toprule
\textbf{Viết tắt} & \textbf{Nguyên văn} & \textbf{Nghĩa dùng trong luận văn} \\
\midrule
\endhead
\bottomrule
\endfoot
PINN & physics-informed neural network & mạng nơ-ron thông tin vật lý \\
FNN & feedforward neural network & mạng nơ-ron truyền thẳng \\
MLP & multilayer perceptron & perceptron nhiều lớp \\
AD & automatic differentiation & vi phân tự động \\
JVP & Jacobian--vector product & tích Jacobi--vectơ (chế độ thuận) \\
VJP & vector--Jacobian product & tích vectơ--Jacobi (chế độ ngược) \\
MSE & mean squared error & sai số bình phương trung bình \\
SGD & stochastic gradient descent & hạ gradient ngẫu nhiên \\
BFGS & Broyden--Fletcher--Goldfarb--Shanno & phương pháp tựa Newton \\
L-BFGS & limited-memory BFGS & BFGS bộ nhớ hạn chế \\
ReLU & rectified linear unit & hàm kích hoạt $\max\{0,z\}$ \\
SiLU & sigmoid linear unit & hàm kích hoạt Swish, $z\,\sigma_{\mathrm{s}}(z)$ \\
LHS & Latin hypercube sampling & lấy mẫu siêu khối Latin \\
RAR & residual-based adaptive refinement & lấy mẫu thích nghi theo phần dư \\
NTK & neural tangent kernel & nhân tiếp tuyến nơ-ron \\
RK4 & Runge--Kutta bậc bốn & lược đồ tích phân thời gian tường minh \\
cPINN & conservative PINN & PINN phân rã miền dạng bảo toàn \\
XPINN & extended PINN & PINN phân rã miền tổng quát \\
\bottomrule
\end{longtable}

\begin{center}
\begin{minipage}{0.92\textwidth}
\small
\textbf{Về hai chữ viết tắt không dùng.} Luận văn viết đầy đủ ``phương trình đạo
hàm riêng'' và ``phương trình vi phân thường'' thay vì PDE và ODE, để tránh lẫn
với cách viết tắt tiếng Việt; riêng ký hiệu $\mathcal{L}$ luôn chỉ toán tử vi
phân và không bao giờ chỉ hàm mất mát (xem quy ước ở đầu
Chương~\ref{chap:pinn}).
\end{minipage}
\end{center}
